\subsection{The exact additional inner freedom at the cusp}

For a general positive $\tau$-fixed quotient action
$\alpha_U$, changing it to $\operatorname{Inn}(h)\alpha_U$
with $h\in C_U(\tau)$ changes the lifted cusp defect
by the continuous homomorphism
\[
 \kappa_\tau:C_U(\tau)\longrightarrow M,
 \qquad h\longmapsto
             \widetilde h\tau\widetilde h^{-1}\tau^{-1}.        \tag{B.S19}
\]
Centrality of $M$ proves independence of the lift and the
homomorphism law, exactly as for (B.S11).  The $S_3$
residue of $h$ belongs to the centralizer of the cusp
transposition, so multiplying $h$ by a suitable power
of $\tau$ makes it pure.  This does not change its
value under (B.S19).  A pure inner change descends
through the first-coordinate map $J\to B$ to the
corresponding inner target change, preserving the
normalized target outer parameter.  Thus this freedom
must be retained in a general cusp-defect test.

\begin{result}{B.2.18}{Proved}{The actual cusp centralizer has no extra central correction}
For the specific universal six-coordinate quotient $U$,
\[
                   C_J(\tau)=\overline{\langle x\rangle},
       \qquad C_U(\tau)=\overline{\langle\tau\rangle},
       \qquad \kappa_\tau=1.                                 \tag{B.S20}
\]
Consequently the additional inner freedom (B.S19)
does not enlarge the defect coset (B.S11) for this $U$.
\end{result}

\begin{proof}
Compose the two actual Schreier actions used in
(B.R18).  The action of the source $\tau=sr$ on its
six pure coordinates is exactly
\[
 (w_1,\ldots,w_6)\longmapsto
 (w_5,w_3,\tau^{-1}w_2\tau,
                    sw_6s,sw_1s,\tau w_4\tau^{-1}).            \tag{B.S21}
\]
In this display the $s,r,\tau$ in the individual
coordinates belong to the target copy of $B$.
The fixed tuples therefore have the form
\[
 (s^e,\tau^t,\tau^t,r^l,s^e,\tau r^l\tau^{-1}),
 \qquad e\in\mathbb Z/2,\quad l\in\mathbb Z/3,
                       \quad t\in\widehat{\mathbb Z}.         \tag{B.S22}
\]
Here we used the actual target centralizers
$C_B(s)=\langle s\rangle$, $C_B(r)=\langle r\rangle$,
and $C_B(\tau)=\overline{\langle\tau\rangle}$.
For the last assertion, a centralizer element reduces
in $S_3$ to a power of the cusp transposition.  Remove
that power of $\tau$ to make it pure.  It centralizes
$\tau^2=x^{-1}$, and the primitive free-group cusp
centralizer is $\overline{\langle x\rangle}$.  This
proves the assertion with the actual pure subgroup;
no centralizer theorem for an arbitrary profinite
word is being used.

It remains to impose membership in the actual $J$.
By (B.R4) and (B.R7), it is enough to reduce (B.S22)
in $\mathcal T^6$ and test its membership in
$J_{\mathcal T}$.  The named target $\tau$ has order
$12$ there.  The exact script
\path{check_universal_cusp_centralizer.g} tests all
$2\cdot12\cdot3=72$ tuples.  Exactly $12$ belong
to $J_{\mathcal T}$; they are precisely
\[
                    e=t\pmod2,\qquad l=-t\pmod3.              \tag{B.S23}
\]
The output records all twelve triples and checks
these two congruences for every one of them.

For every full profinite $t$ satisfying these conditions,
the tuple (B.S22) equals the literal tuple $\eta(x^t)$
from (B.R10): its last coordinate is
$sr^{-t}s=\tau r^{-t}\tau^{-1}$.  Thus every pure
centralizer element is the corresponding power of
the source $x$, proving the first assertion of
(B.S20).  Removing the possible cusp-transposition
residue from any element of $C_U(\tau)$ then proves
the second assertion.  The named $\tau$ has infinite
procyclic order in $U$, as witnessed by the infinite
target $\tau$ coordinate of $x=\tau^{-2}$.

Finally every element of this centralizer has the
actual lift $\tau^t$ in $U^{\mathrm c}$, which
commutes with $\tau$.  Equation (B.S19) is independent
of the lift, so it is identically one.  This proves
the final assertion of (B.S20).
\end{proof}
